% ============================================================================
%  Part II (training) -- Section 6: The training algorithm
%  \input-ready fragment. Uses only macros from preamble.tex.
%  Cross-part refs (defined elsewhere): eq:elbo, eq:predmoments, eq:fbar
%  (Part II inference).
% ============================================================================

\section{The training algorithm: EM-style coordinate ascent}\label{sec:training}

The model assembled in Parts~I--II has three groups of free parameters, and all
three sit on the \emph{single} evidence lower bound $\ELBO$ of
\eqref{eq:elbo}:

\begin{enumerate}[label=(\alph*),leftmargin=2.2em,itemsep=1pt]
  \item the \textbf{variational moments} $\mb,\Vb$ of the eigenspace posterior
        $q(\latt)=\Gauss(\vm,\vV)$;
  \item the \textbf{rate parameters} of the Poisson likelihood, the gain
        $\gain$ and the offset $\latz$, defining
        $\rate=\exp(\gain\lat+\latz)$;
  \item the \textbf{kernel hyperparameters}
        $\{\sigz,\Amp,\bwid,\smoo,\varepsilon_{0x},\varepsilon_{0y}\}$ of the
        structured arc-cosine prior (Part~I).
\end{enumerate}

Training maximises the one bound $\ELBO$ by \textbf{block coordinate ascent}:
each step improves one block while holding the other two fixed. The three steps
are the E-step (variational moments), the F-step (rate parameters), and the
M-step (kernel hyperparameters), run once per outer iteration in a fixed order
(\S\ref{subsec:training-overview}).

% ---------------------------------------------------------------------------
\subsection{Overview: three coordinate blocks, one ELBO, one loop}
\label{subsec:training-overview}

Each step optimises one block and freezes the rest:

\begin{center}\small
\begin{tabular}{@{}l l l >{\raggedright\arraybackslash}p{2.35in}@{}}
\toprule
Step & Free block & Held fixed & Method \\
\midrule
E & $\mb,\Vb$ & kernel; $\gain,\latz$; $\eb$ &
    closed-form Newton (no automatic differentiation) \\
F & $\gain$ ($\latz$ closed-form) & $\mb,\Vb$ (via $\pmean,\pvar$); kernel &
    LBFGS on $\log\gain$, or damped Newton \\
M & $\{\sigz,\Amp,\bwid,\smoo,\rfc\}$ & $\mb,\Vb$; $\gain,\latz$; $\eb$ &
    LBFGS $+$ gradients (analytical or automatic) \\
\bottomrule
\end{tabular}
\end{center}

\noindent(In the M-row, the centre $\rfc=(\varepsilon_{0x},\varepsilon_{0y})$ of
the localized region counts as the two coordinate hyperparameters, so the free
block has six scalars.)

\paragraph{Loop order (one outer iteration).}
Per outer iteration:

\begin{enumerate}[leftmargin=2.4em,itemsep=1pt]
  \item \textbf{Re-sync the eigenbasis.} Rebuild the eigenbasis $\eb$ from the
        hyperparameters the \emph{previous} M-step moved, then refresh
        $\pmean,\pvar,\fbar$ --- this is where the M-step's deferred $\eb$ update
        lands.
  \item \textbf{E-step} --- one or more closed-form Newton steps on
        $\mb,\Vb$ (\S\ref{subsec:estep}).
  \item \textbf{F-step} --- fit $\gain$ (and closed-form $\latz$), unless the
        F-step is interleaved into the E-loop (\S\ref{subsec:fstep}).
  \item \textbf{Metrics $+$ ELBO $+$ early stopping} --- computed
        \emph{deliberately before} the M-step.
  \item \textbf{M-step} --- update the kernel hyperparameters
        (\S\ref{subsec:mstep}); the $\eb$ update is \emph{deferred} to step~1 of
        the next iteration.
\end{enumerate}

\begin{remark}[Why metrics precede the M-step]
The M-step changes the kernel parameters \emph{without} recomputing the
eigenbasis $\eb$ (that is deferred, for cost and stability). Every logged
quantity --- the ELBO, the early-stopping metric --- must therefore be read while
the model state is still internally consistent, i.e.\ \emph{before} the M-step
perturbs the kernel. Early stopping is on the ELBO, with patience and
best-restore.
\end{remark}

% ---------------------------------------------------------------------------
\subsection{The E-step: a Newton update of the variational posterior}
\label{subsec:estep}

At \emph{fixed} kernel and fixed $\gain,\latz$, the E-step maximises the ELBO
\eqref{eq:elbo} over the variational moments. Writing the inducing-space
posterior $q(\latt)=\Gauss(\vm,\vV)$ against the prior $p(\latt)=\Gauss(0,\Ktil)$,

\begin{equation}
\ELBO(\vm,\vV)
  = \underbrace{\sum_i \E_{q(\lat_i)}\!\big[\log\Poi(\spk_i\mid\lat_i)\big]}_{\ELBO_{\mathrm{data}}}
    \;-\; \KLdiv\!\big(q(\latt)\,\|\,p(\latt)\big),
\label{eq:estep-obj}
\end{equation}
with the projected posterior moments $\pmean_i,\pvar_i$ of \eqref{eq:predmoments}
and the KL contributing, through the E-step variables,
$-\tfrac12\vm\transpose\Ktil^{-1}\vm$ (couples to $\vm$) and
$\tfrac12\log|\vV|-\tfrac12\Tr(\Ktil^{-1}\vV)$ (couples to $\vV$).

\paragraph{Poisson expected log-likelihood and its moment derivatives.}
The data term is closed-form via the Gaussian moment-generating function, with
the expected rate $\fbar_i$ of \eqref{eq:fbar}:
\begin{equation}
\ELBO_{\mathrm{data}}
  = \sum_i\big[\,\spk_i(\gain\pmean_i+\latz)-\fbar_i\,\big],
\qquad
\fbar_i=\exp\!\big(\gain\pmean_i+\tfrac12\gain^2\pvar_i+\latz\big),
\end{equation}
(the $-\log\spk_i!$ constant is dropped). The three point-derivatives that
assemble the Newton step follow by the chain rule on $\fbar_i$:
\begin{equation}
\frac{\partial\fbar_i}{\partial\pmean_i}=\gain\,\fbar_i,\qquad
\frac{\partial\fbar_i}{\partial\pvar_i}=\tfrac12\gain^2\,\fbar_i,\qquad
\frac{\partial}{\partial\pmean_i}\,\E_{q}[\log\Poi]=\gain(\spk_i-\fbar_i).
\end{equation}

\paragraph{The two Newton helpers.}
Accumulating those derivatives through the projection $\pmean_i=\kvec_i\transpose\Ktil^{-1}\vm$
defines the gradient helper $\gE$ and the (Poisson-Fisher) curvature $\GE$:
\begin{equation}
\gE \;=\; \gain\sum_i \kvec_i(\spk_i-\fbar_i)\ \in\Reals^{\Mind},
\qquad
\GE \;=\; \gain^2\sum_i \fbar_i\,\kvec_i\kvec_i\transpose\ \in\Reals^{\Mind\times\Mind}.
\label{eq:estep-gG}
\end{equation}
$\gE$ is the kernel-projected mismatch between observed and expected counts;
$\GE$ is the Poisson Fisher information and is PSD because every $\fbar_i>0$.
(Note $\gE$ carries the subscript $\mathrm{E}$ specifically to keep it distinct
from the log-rate $\lgr=\gain\lat+\latz$ of the utility part.)

\paragraph{The $\vV$-update (closed form).}
With $\pmean_i$ independent of $\vV$, setting $\nabla_{\vV}\ELBO=0$ gives
$\vV^{-1}=\Ktil^{-1}(\Ktil+\GE)\Ktil^{-1}$, hence
\begin{equation}
\boxed{\;\vV \;=\; \Ktil\,(\Ktil+\GE)^{-1}\,\Ktil\;}
\label{eq:estep-V}
\end{equation}
--- a \emph{symmetric} sandwich of two symmetric-PD factors, so $\vV$ is a valid
covariance. This symmetry is the substance of the ``corrected'' E-step: the
superseded form $\vV=(\Ktil+\GE)^{-1}\Ktil$ is a product of two non-commuting
symmetric matrices and is therefore \emph{not} symmetric (Cholesky/PSD failures).

\paragraph{The $\vm$-update.}
The exact Newton step $\vm\leftarrow\vm-H_{\vm}^{-1}\nabla_{\vm}\ELBO$, with
$\nabla_{\vm}\ELBO=\Ktil^{-1}(\gE-\vm)$ and Hessian
$H_{\vm}=-\Ktil^{-1}(\Ktil+\GE)\Ktil^{-1}$, gives the corrected form
$\vm\leftarrow\vm+\Ktil(\Ktil+\GE)^{-1}(\gE-\vm)$. The map actually used is
\emph{not} this; it is the algebraically older
\begin{equation}
\boxed{\;\vm \;\leftarrow\; \Ktil\,(\Ktil+\GE)^{-1}\big(\GE\,\Ktil^{-1}\vm+\gE\big)\;}
\label{eq:estep-m}
\end{equation}
(equal to $\vV_{\mathrm{new}}(\GE\Ktil^{-1}\vm+\gE)$ with $\vV_{\mathrm{new}}$ from
\eqref{eq:estep-V}). The two maps differ only in the first term
($\Ktil(\Ktil+\GE)^{-1}\GE\Ktil^{-1}\vm$ vs.\ $\GE(\Ktil+\GE)^{-1}\vm$); they
coincide iff $\Ktil$ and $\GE$ commute.

\caveat{The $\vm$-update map \eqref{eq:estep-m} differs \emph{per iteration} from
the corrected Newton step whenever $\Ktil$ and $\GE$ do not commute, but both
maps share the same fixed point $\vm^\star=\gE$, so the \emph{converged}
posterior is identical --- the discrepancy is a transient path difference, not a
different solution. It is left as-is (it converges in practice). The $\vV$-update
\eqref{eq:estep-V} is unambiguously correct, and it is $\vV$ that the
active-learning variance estimates depend on.}

\begin{remark}[Fixed-point equivalence]
Let $T_{\mathrm N}(\vm)=\vm+\Ktil(\Ktil+\GE)^{-1}(\gE-\vm)$ (corrected) and
$T_{\mathrm C}(\vm)=\Ktil(\Ktil+\GE)^{-1}(\GE\Ktil^{-1}\vm+\gE)$ (older map). At
$\vm=\gE$:
$T_{\mathrm N}(\gE)=\gE$, and
$T_{\mathrm C}(\gE)=\Ktil(\Ktil+\GE)^{-1}(\GE\Ktil^{-1}+I)\gE
   =\Ktil(\Ktil+\GE)^{-1}(\GE+\Ktil)\Ktil^{-1}\gE=\Ktil\Ktil^{-1}\gE=\gE$.
Both are valid fixed-point iterations onto the same stationary condition
$\vm^\star=\gE=\gain\sum_i\kvec_i(\spk_i-\fbar_i)$.
\end{remark}

\subsubsection{Why the update is cheap: the eigenspace realisation}
\label{subsubsec:estep-eigen}

The E-step never touches the $\Mind$-dimensional inducing space directly. It works
in the eigenbasis $\eb$ of $\Ktil$, keeping only the $\nb$ leading directions
(eigenvalues above $\max(\lambda_{\max}\!\cdot\!\mathrm{tol},\,\mathrm{tol})$,
$\mathrm{tol}=10^{-4}$; in practice $\nb\approx10\text{--}11$ while $\Mind$ may be
large). With $\Ktil=\eb\evals\eb\transpose$ the projected gram
$\Ktilb=\eb\transpose\Ktil\eb=\diago(\evals)$ is \emph{diagonal}, so $\Ktil^{-1}$
is trivial. Using the eigenspace projection $\amat=K\Ktil^{-1}$
(shape $\Ntr\times\nb$), one forms the \emph{transformed} helpers:
\begin{align}
g_{\mathrm E,b} &= \gain\,\amat\transpose(\spk-\fbar)\;=\;\Ktil^{-1}\gE
   &&(\text{eigenbasis coords}),\\
G_{\mathrm E,b} &= \gain^2\,\amat\transpose\big(\fbar\odot\amat\big)\;=\;\Ktil^{-1}\GE\Ktil^{-1},
\end{align}
and applies \eqref{eq:estep-V}--\eqref{eq:estep-m} entirely in the tiny
$\nb$-space:
\begin{equation}
\Vb \leftarrow \big(I+\Ktilb\,G_{\mathrm E,b}\big)^{-1}\Ktilb,\qquad
\mb \leftarrow \Vb\big(G_{\mathrm E,b}\,\mb+g_{\mathrm E,b}\big),\qquad
\Vb \leftarrow \tfrac12(\Vb+\Vb\transpose).
\end{equation}
The first line equals \eqref{eq:estep-V} exactly, via the identity
$(I+PQ^{-1})^{-1}=Q(Q+P)^{-1}$ with $P=\GE,Q=\Ktil$:
$(I+\Ktilb G_{\mathrm E,b})^{-1}\Ktilb=(I+\GE\Ktil^{-1})^{-1}\Ktil
=\Ktil(\Ktil+\GE)^{-1}\Ktil$. Prior initialisation is $\mb=0$,
$\Vb=\Ktilb=\evals$. The dimensional collapse $\Mind\!\to\!\nb\approx10$ ---
not any special-casing --- is what makes the Newton step cheap.

\subsubsection{Iteration, damping, and the divergence guard}

The E-step performs a \emph{single, undamped, closed-form} Newton step; iteration
and stabilisation are handled around it:

\begin{itemize}[leftmargin=1.4em,itemsep=1pt]
  \item \textbf{Coupling variable $\fbar$.} After each step the posterior is
        recomputed and $\fbar=\exp(\gain\pmean+\tfrac12\gain^2\pvar+\latz)$
        refreshed, then fed into the next step's $g_{\mathrm E,b},G_{\mathrm E,b}$.
        $\gain,\latz$ are held fixed during the E-step.
  \item \textbf{Revert-on-divergence guard (the effective damping).} Each iteration
        snapshots $(\mb,\Vb)$; if the post-step $\fbar$ mean or max exceeds a
        threshold, or $\fbar$ has NaNs, it \emph{reverts} $(\mb,\Vb)$ (and
        $\gain,\latz$ if interleaved), recomputes $\fbar$, and breaks the loop.
  \item \textbf{Optional interleaved F-step.} With interleaving enabled, a
        \emph{damped} Newton update of $(\gain,\latz)$ (\S\ref{subsec:fstep},
        $\alpha=0.25$) runs after every E-iteration, so the rate parameters
        track $(\mb,\Vb)$ as they move.
\end{itemize}

% ---------------------------------------------------------------------------
\subsection{The F-step: fitting the rate function}
\label{subsec:fstep}

Here ``F'' denotes the \textbf{rate function} $\rate$: the F-step fits the
Poisson-rate parameters $\gain,\latz$ that define
$\fbar=\exp(\gain\pmean+\tfrac12\gain^2\pvar+\latz)$, holding the variational
moments (through frozen $\pmean,\pvar$) and the kernel fixed. Here $\gain$ is the
only searched variable; $\latz$ is solved in closed form for the current $\gain$.

\paragraph{Closed-form $\latz$ given $\gain$.}
From $\partial\ELBO/\partial\latz=0$ on
$\sum_i[\spk_i\latz-\exp(\latz)\exp(\gain\pmean_i+\tfrac12\gain^2\pvar_i)]$:
\begin{equation}
\boxed{\;
\latz \;=\; \log\!\Big(\textstyle\sum_i\spk_i\Big)
       \;-\; \log\!\Big(\textstyle\sum_i\exp\!\big(\gain\pmean_i+\tfrac12\gain^2\pvar_i\big)\Big)
\;}
\label{eq:fstep}
\end{equation}
(requires $\sum_i\spk_i>0$; a zero-count target is a data problem). It is
re-evaluated whenever $\gain$ changes.

\paragraph{LBFGS on $\gain$ with an analytic gradient.}
The F-step runs LBFGS over $\log\gain$, re-solving $\latz$ at each evaluation and
using the closed-form derivative, with $\pmean,\pvar$ the frozen inputs:
\begin{equation}
\frac{\partial\ELBO}{\partial\gain}
  = \sum_i\big[\spk_i\pmean_i-(\pmean_i+\gain\pvar_i)\,\fbar_i\big],
\qquad
\frac{\partial\ELBO}{\partial(\log\gain)}=\gain\frac{\partial\ELBO}{\partial\gain}.
\end{equation}
Since LBFGS minimises $-\ELBO$, the gradient it consumes with respect to
$\log\gain$ is $-\gain\,\partial\ELBO/\partial\gain$. Guards reject a trial step
whose $\latz$ over/underflows or whose $\fbar$ exceeds the mean/max thresholds,
and revert $(\gain,\latz)$ if the final $\latz$ overflows.

\paragraph{Interleaved joint damped Newton on $(\gain,\latz)$.}
With interleaving enabled, a joint damped Newton step on $\psi=[\gain,\latz]$ runs
\emph{inside} the E-loop. With per-input
$\rate_i=\exp(\gain\pmean_i+\tfrac12\gain^2\pvar_i+\latz)$ and
$d_i=\pmean_i+\gain\pvar_i$:
\begin{equation}
R=\begin{bmatrix}\sum_i(\spk_i\pmean_i-d_i\rate_i)\\[2pt]\sum_i(\spk_i-\rate_i)\end{bmatrix},
\quad
H=-\begin{bmatrix}\sum_i(\pvar_i\rate_i+d_i^2\rate_i)&\sum_i d_i\rate_i\\[2pt]
                  \sum_i d_i\rate_i&\sum_i\rate_i\end{bmatrix},
\quad
\psi\leftarrow\psi-\alpha\,H^{-1}R,\ \ \alpha=0.25.
\end{equation}
$H$ is negative-definite (the log-likelihood is concave in $\psi$); the damping
$\alpha=0.25$ and the same $\fbar$-threshold rejections stabilise the step.

% ---------------------------------------------------------------------------
\subsection{The M-step: updating the kernel hyperparameters}
\label{subsec:mstep}

The M-step maximises the ELBO \eqref{eq:elbo} over the kernel/prior
hyperparameters while $\mb,\Vb$, the rate parameters $\gain,\latz$, and
\emph{the eigenbasis $\eb$} are held fixed. The free hyperparameters (all owned
by the arc-cosine prior) are the six
\begin{equation}
\theta\in\{\ \sigz,\ \Amp,\ \bwid,\ \smoo,\ \varepsilon_{0x},\ \varepsilon_{0y}\ \}.
\end{equation}
The eigenbasis $\eb$ is held fixed as an approximation (for cost and stability)
and re-synced at the start of the next outer iteration. This block is solved by
gradient-based optimisation (e.g.\ LBFGS on the raw, unconstrained parameters);
the hyperparameter gradients $\partial_\theta\ELBO$ are obtained by automatic or
analytical differentiation of the eigenspace ELBO.

\begin{remark}[The amplitude $\Amp$ is optimized by default]
The amplitude prefactor $\Amp$ of the structured prior
$\Cmat=\Amp\,\loc\,\Csm\,\loc\transpose$ (Part~I) is a genuine sixth kernel
hyperparameter, not a fixed constant, and is optimized alongside the other five
by default. The optional flag \texttt{fix\_Amp} (default off) instead pins
$\Amp=1$, recovering the amplitude-free model; left at its default, $\Amp$ trains
with the rest.
\end{remark}
